It is natural to consider Tate cohomology of~$V$ for elements of composite order, but there are no results in the literature. Tate cohomology does not yield vector spaces over a field when the element has composite order, but Brauer characters are only defined for such objects. Specifically, Tate cohomology yields torsion modules over the ground ring. We can't use the Brauer character for such modules because it is defined only for representations defined over a~field and $p$-regular elements. So, we need to define a generalized Brauer character.

In Section~\ref{section2}, we introduce some terminology in this paper. In Section~\ref{section3}, we summarize the modular moonshine conjecture for elements of prime order. In Section~\ref{section4}, we consider ramified extensions $R_p$ of the $p$-adic integers $\mathbb{Z}_p$ and calculate Tate cohomology. In Section~\ref{section5}, to generalize the modular moonshine conjectures, we define a generalized Brauer character and show that the generalized super Brauer character of Tate cohomology is a linear combination of traces. This is a generalization of \cite[Proposition~2.2]{BR} to the composite order case. Using this result, we give a counterexample to Borcherds's conjecture about vanishing of Tate cohomology. We propose a~weaker conjecture about vanishing of Tate cohomology.
